\section{Auswertung}
\label{sec:Auswertung}


\subsection{BB84 Schlüsselgenerierung und Nachrichtenverschlüsselung}
\label{sec:a1}
In diesem Abschnitt wird sich zuallererst mit der Funktionsweise der Quantenkryptographie vertraut gemacht. 
Zur Schlüsselgenerierung wurden 52 Bits von Alice an Bob übertragen. Bob hat jeweils eine zufällige der beiden Basen gewählt ($\times$ oder $+$) 
und die 52 Bits von Alice empfangen. Anschließend haben Alice und Bob ihre Basen über einen öffentlichen Kanal verglichen und die Bits, bei denen sie die gleiche Basis gewählt haben, behalten.
In diesem Fall waren es 27 Bits, die genutzt werden konnten. Zur Verifikation der Funktionsfähigkeit werden hier die Zahlenfolgen der 27 Bits, also die 27 Nullen oder Einsen, die von Alice geschickt und von Bob 
empfangen wurden, verglichen.
\begin{align*}
  \text{Alice:} & \quad 1 \, 1 \, 1 \, 0 \, 1 \, 0 \, 1 \, 0 \, 0 \, 1 \, 1 \, 0 \, 1 \, 0 \, 1 \, 1 \, 0 \, 1 \, 0 \, 0 \, 0 \, 1 \, 1 \, 0 \, 1 \, 0 \, 1 \\
  \text{Bob}: & \quad 1 \, 1 \, 1 \, 0 \, 1 \, 0 \, 1 \, 0 \, 0 \, 1 \, 1 \, 0 \, 1 \, 0 \, 1 \, 1 \, 0 \, 1 \, 0 \, 0 \, 0 \, 1 \, 1 \, 0 \, 1 \, 0 \, 1
\end{align*}

\noindent Da die Zahlenfolgen übereinstimmen, haben somit Alice und Bob jetzt den gleichen Schlüssel. Das ergibt eine Fehlerrate von $\frac{N_\text{Error}}{N_\text{shifted Bits}} = 0$, da alle Bits übereinstimmen.
Anschließend wird mithilfe des Schlüssels von Alice eine Nachricht verschlüsselt.
Die Nachricht lautet: ``ALSO''. Diese Nachricht wird in Binärcode umgewandelt und 
mithilfe des ersten 20 Bits des Schlüssels von Alice verschlüsselt.
Beim Verschlüsseln wird der Schlüssel über die Nachricht gelegt und in der verschlüsselten
Nachricht wird bei gleichen Zahlen, also (1 und 1) oder (0 und 0), eine Null als
verschlüsseltes Bit vermerkt und bei ungleichen Bits (1 und 0) oder (0 und 1) eine
1 als verschlüsseltes Bit vermerkt (siehe \autoref{sec:bb84}). Das Prinzip ist das gleiche wie bei einem 
XOR-Gatter in der Digitaltechnik.
\begin{align*}
\text{ALSO} & \quad \rightarrow \quad 00000 \, 01011 \, 10010 \, 01110 \\
\text{Schlüssel} & \quad \rightarrow \quad 11101 \, 01001 \, 10101 \, 10100\\
\text{Verschlüsselte Nachricht} & \quad \rightarrow \quad 11101 \, 00010 \, 00111 \, 11010
\end{align*}
Diese verschlüsselte Nachricht kann nun öffentlich an Bob gesendet werden und Bob kann durch Anwenden des gleichen Prinzips auf den Binärcode der Buchstaben schließen und die Nachricht entschlüsseln.
\begin{align*}
\text{verschlüsselte Nachricht} & \quad \rightarrow \quad 11101 \, 00010 \, 00111 \, 11010 \\
\text{Schlüssel} & \quad \rightarrow \quad 11101 \, 01001 \, 10101 \, 10100\\
\text{Entschlüsselte Nachricht} & \quad \rightarrow \quad 00000 \, 01011 \, 10010 \, 01110 
\end{align*}

\subsection{Abhörversuch von Eve}
\label{sec:a2}
In diesem Abschnitt wird der Abhörversuch von Eve beschrieben. Eve versucht die Nachricht von Alice
abzufangen und zu entschlüsseln. Dazu wählt Eve zufällig eine Basis ($\times$ oder $+$) und misst die 52 Bits von Alice. 
Damit Bob keinen Verdacht schöpft und überhaupt etwas empfängt, muss Eve die gemessenen Bits an Bob weiterleiten. Da Eve aber nicht weiß, 
welche Basis Alice gewählt hat, kann es passieren, dass sie die falsche Basis wählt und somit falsche Bits an Bob weiterleitet.
Falls Bob dann die falsche Basis wählt, verfallen die Bits beim Vergleichen der Basen. Wenn Bob jedoch die richtige wählt, nachdem Eve die falsche gewählt hat, kann es 
dazu kommen, dass die von Eve in die falsche Basis gedrehten und von Bob dann wieder in die richtige Basis gedrehten Bits falsch sind. Also hat Bob die gleiche Basis wie Alice 
gewählt, aber die falschen Bits von Eve erhalten. Das führt zu einer Fehlerrate von $\frac{N_\text{Error}}{N_\text{shifted Bits}} > 0$. 

\noindent Es wurden wieder 52 Bits losgeschickt, dann von Eve abgehört und das Ergebnis an Bob weitergeleitet. Die jeweiligen 
Tabellen zu den Messreihen sind im Anhang zu finden. Im Folgenden werden die Bits verglichen, bei denen Alice und Bob die gleiche 
Basis gewählt haben. Daraus wird dann die Fehlerrate bestimmt, die auf Eves Zwischenschalten zurückzuführen ist.
\begin{table}[H]
\centering
\begin{tabular}{c|c|c|c|c|c}
\textbf{Alice} &  \textbf{Bob} & Übereinstimmung &\textbf{Alice} &  \textbf{Bob} & Übereinstimmung \\ \hline
0 & 0 &  &1 & 1 &  \\
0 & 0 &  &0 & 0 &  \\
1 & 1 &  &1 & 1 &  \\
1 & 1 &  &1 & 0 & x \\
0 & 0 &  &0 & 0 &  \\
1 & 0 & x &1 & 1 &  \\
0 & 1 & x &1 & 1 &  \\
1 & 1 &  &0 & 0 &  \\
0 & 0 &  &1 & 0 & x \\
0 & 0 &  &1 & 1 &  \\
1 & 1 &  &0 & 1 & x\\
1 & 1 &  &0 & 0 &  \\
0 & 1 & x &1 & 1 &  \\
1 & 0 & x &0 & 0 &  \\
1 & 1 &  &1 & 1 &  \\
0 & 1 & x &1 & 1 &  \\
\hline
\end{tabular}
\caption{Vergleich der shifted Bits von Alice und Bob bei einem Abhörversuch von Eve}
\label{tab:vergleichEve}
\end{table}
\noindent Die gesendeten und empfangenen Bits, welche durch eine gleiche Basis von
Alice und Bob ausgewählt wurden, sind \autoref{tab:vergleichEve} zu entnehmen. Die
fehlerhaften Bits sind mit einem Kreuz markiert und belaufen sich auf acht Fehler
bei 32 ausgewählten verglichenen Bits. Das führt zu einer Fehlerrate von 
\begin{equation}
 \frac{N_\text{Error}}{N_\text{shifted Bits}} = \frac{8}{32} = 0.25 
\end{equation}
Das bedeutet, dass 25\% der verglichenen Bits fehlerhaft sind, was auf Eves Abhörversuch zurückzuführen ist.
% Eve und Bob können nun diese Fehlerquote als Indiz für Eves Abhörversuch nutzen. Dafür würden sie in der Praxis einen Teil ihres Schlüssels vergleichen, um Fehler aufzudecken. 
% Falls Fehler zu verzeichnen sind, ist Eve aufgedeckt und der Schlüssel sollte von Alice und Bob nicht verwendet werden. 
% Hier wurden alle Bits des Schlüssels verglichen, da es sonst zu wenig Vergleichsmaterial gegeben hätte.


\subsection{Schlüsselgenerierung mit Decoy States}
\label{sec:a3}
In diesem Abschnitt wird die Schlüsselgenerierung mit Decoy States beschrieben. 
Es werden statt Bits jetzt Trits gesendet (0, 1 oder 2) und das wieder entweder in
der $\times$ Basis oder in der $+$ Basis.
Die Kodierungen sind wie folgt in \autoref{tab:ternary_protocol} dargestellt. 

\definecolor{colorRed}{RGB}{200, 0, 0}
\definecolor{colorGreen}{RGB}{0, 150, 0}

\begin{table}[H]
\centering
\caption{Kodierung der Trits in den beiden Basen und den Decoy States}
\label{tab:ternary_protocol}
\renewcommand{\arraystretch}{1.3} % Erhöht den Zeilenabstand für bessere Lesbarkeit
\begin{tabular}{|c|c|c|c|}
\hline
\text{Basis} & \text{Trit} & \text{Polarisationsrichtung} &\text{decoy state} \\ \hline
 & \textbf{0} & \textcolor{colorRed}{$0^\circ$} &  \\ \cline{2-3}
\textbf{[+]} & \textbf{1} & \textcolor{colorRed}{$90^\circ$} &  \\ \cline{2-3}
 & \textbf{2} & \textcolor{colorGreen}{$0^\circ$} &  \\ \cline{2-4} 
 & \multicolumn{2}{c|}{} & \textcolor{colorGreen}{$90^\circ$} \\ \hline
 & \textbf{0} & \textcolor{colorGreen}{$-45^\circ$} &  \\ \cline{2-3}
\textbf{[x]} & \textbf{1} & \textcolor{colorGreen}{$45^\circ$} &  \\ \cline{2-3}
 & \textbf{2} & \textcolor{colorRed}{$-45^\circ$} &  \\ \cline{2-4} 
 & \multicolumn{2}{c|}{} & \textcolor{colorRed}{$45^\circ$} \\ \hline
\end{tabular}
\end{table}

\noindent Bob empfängt diese ausgesendeten Bits und wandelt sie wieder in Trits um. 
Diese Umwandlung erfolgt quasi analog zu Alices Kodierung der Trits, nur müssen 
jetzt entweder die roten oder grünen Bits in der jeweiligen Phase interpretiert 
werden.

% Definition der Farben aus dem Praktikum
\definecolor{colorRed}{RGB}{200, 0, 0}
\definecolor{colorGreen}{RGB}{0, 150, 0}

\begin{table}[H]
\centering
\caption{Umwandlung der Signalzustände in Trits bei Bob (Basis, Trit, Farbe, Polarisation)}
\label{tab:signal_states_mapping}
\renewcommand{\arraystretch}{1.2}
\begin{tabular}{cccccc}
\toprule
\text{Basis} & \text{Trit} & \text{Bezeichnung} & \text{Laser Farbe} & \text{Polarisationswinkel} \\ 
\midrule
 & 0 & \texttt{r0+} & \textcolor{colorRed}{Rot} & $0^\circ$ \\ \cmidrule{2-5}
\textbf{[+]} & 1 & \texttt{r1+} & \textcolor{colorRed}{Rot} & $90^\circ$ \\ \cmidrule{2-5}
 & 2 & \texttt{g0+} & \textcolor{colorGreen}{Grün} & $0^\circ$ \\ 
\midrule
 & 0 & \texttt{g0x} & \textcolor{colorGreen}{Grün} & $-45^\circ$ \\ \cmidrule{2-5}
\textbf{[x]} & 1 & \texttt{g1x} & \textcolor{colorGreen}{Grün} & $45^\circ$ \\ \cmidrule{2-5}
 & 2 & \texttt{r0x} & \textcolor{colorRed}{Rot} & $-45^\circ$ \\ 
\bottomrule
\end{tabular}
\end{table}

\noindent Nun werden 52 Trits in zufälligen Basen verschickt und es wird wie in
\autoref{sec:a1} verfahren. Nach Verwerfen der Trits, bei denen Alice und Bob
unterschiedliche Basen gewählt haben, bleiben 32 Trits übrig. Diese werden in 
\autoref{eqn:trits} verglichen.
\begin{align*}
  \label{eqn:trits}
  \text{Alice:} & \quad 1 \, 2 \, 2 \, 1 \, 1 \, 2 \, 2 \, 1 \, 2 \, 0 \, 1 \, 0 \, 1 \, 0 \, 1 \, 0 \, 0 \, 2 \, 2 \, 0 \, 0 \, 2 \, 0 \, 2 \, 0 \, 2 \, 0 \, 1 \, 1 \, 0 \, 2 \, 2\\
  \text{Bob}: & \quad 1 \, 2 \, 2 \, 1 \, 1 \, 2 \, 2 \, 1 \, 2 \, 0 \, 1 \, 0 \, 1 \, 0 \, 1 \, 0 \, 0 \, 2 \, 2 \, 0 \, 0 \, 2 \, 0 \, 2 \, 0 \, 2 \, 0 \, 1 \, 1 \, 0 \, 2 \, 2
\end{align*}
\noindent Es ist zu erkennen, dass alle Trits von Alice und Bob übereinstimmen, was zu einer Fehlerrate von $\frac{N_\text{Error}}{N_\text{shifted Trits}} = 0$ und auch 
zu einem Anteil an Decoy States von $\frac{{N_{decoy}}}{{N_{shifted Bits}}} = 0$ führt. Diese wird im folgenden Abschnitt mit der Fehlerrate nach Dazuschalten 
von Eve verglichen.

\subsection{Abhörversuch von Eve bei der Schlüsselgenerierung mit Decoy States}
\label{sec:a4}
In diesem Abschnitt wird der Abhörversuch von Eve bei der Schlüsselgenerierung mit
Decoy States beschrieben. Es werden wieder 52 Trits von Alice an Bob gesendet, aber
dieses mal mit Decoy States. Eve versucht die Nachricht von Alice abzufangen und zu
entschlüsseln. Wenn sie die falsche Basis wählt und somit falsche Trits an Bob 
weiterleitet, führt das zu einer Fehlerrate von $\frac{N_\text{Error}}{N_\text{shifted Trits}} > 0$
und zur Entstehung von Decoy States, welche von Alice und Bob als Indiz für Eves
Abhörversuch genutzt werden können. Erhält Bob also so einen Decoy State bei gleicher
Basiswahl wie Alice, ist Eve aufgedeckt und der Schlüssel sollte von Alice und Bob
nicht verwendet werden. Die gesendeten sowie empfangenen Trits sind in \autoref{tab:vergleichEveDecoy}
dargestellt. Die fehlerhaften Trits sind mit einem Kreuz markiert und die Decoy
States als solches betitelt. Die Messreihen sind dem Anhang zu entnehmen.

\begin{table}[H]
\centering
\begin{tabular}{c|c|c|c|c|c}
\textbf{Alice} &  \textbf{Bob} & Übereinstimmung &\textbf{Alice} &  \textbf{Bob} & Übereinstimmung \\ \hline
1 & 2     &  x   & 2 & 0     &  x  \\
0 & Decoy &     & 2 & Decoy &    \\    
2 & 0     &  x   & 0 & 1     &  x  \\
2 & 0     &  x   & 0 & 0     &    \\
1 & Decoy &     & 2 & 0     &  x  \\
1 & 1     &     & 0 & 0     &    \\
2 & Decoy &     & 1 & Decoy &    \\    
2 & 0     &  x   & 2 & 0     &  x  \\
1 & 0     &  x   & 0 & 1     &  x  \\
2 & 0     &  x   & 2 & 0     &  x  \\
0 & 0     &     & 0 & Decoy &    \\    
0 & 0     &     & 1 & 1     &    \\
0 & 0     &     & 1 & 0     &    x\\
1 & 1     &     & 2 & Decoy &    \\    
0 & 0     &     & 0 & 0     &    \\
1 & 1     &     & 2 & 0     &   x \\
0 & 0     &     & 1 & 1     &    \\
0 & 0     &     &   &       &    \\
\hline
\end{tabular}
\caption{Vergleich der shifted Bits von Alice und Bob bei einem Abhörversuch von Eve}
\label{tab:vergleichEveDecoy}
\end{table}

\noindent Damit ergeben sich Error Rate und Decoy Rate bei 33 behaltenen gleichen Basen
und darin 14 Fehlern und 7 Decoys zu 
\begin{align}
  \frac{N_\text{Error}}{N_\text{shifted Bits}} = \frac{14}{33} = 0.42 \\
  \frac{N_\text{Decoy}}{N_\text{shifted Bits}} = \frac{7}{33} = 0.21
\end{align}